%!TEX TS-program = pdflatex
% vim: set fenc=utf-8
% -*- coding: UTF-8; -*-
%!TEX encoding = UTF-8

\documentclass[twoside]{USYDReport}

% =============================================
%                 填写封面信息
% =============================================
\title{Your Report Title Goes Here} % 你的报告/论文题目
\author{Firstname Lastname}         % 你的姓名
\stunum{123456789}                  % 你的学号 (9 digits)
\coursecode{XXXX0000}               % 课程代码 (e.g., IDEA9106)
\tutorial{Tutorial [Number]}        % Tutorial 班级
\groupnum{Group [Number/Name]}      % 小组号或名称 (如果是个人作业可以留空)
\tutor{Your Tutor's Name}           % Tutor 或 Lecturer 姓名

% =============================================
%                 加载额外的宏包
% =============================================
\usepackage{float}      % 控制图表浮动位置 (H)
\usepackage{ragged2e}   % 文本对齐
\usepackage{xspace}     % 智能空格控制
\usepackage{pifont}     % 更多特殊符号
\usepackage{multicol}   % 多栏排版
\usepackage{multirow}   % 表格跨行
\usepackage{hyperref}   % 生成可点击的超链接
\usepackage{newtxmath}  % 数学公式字体匹配 Times
\usepackage{adjustbox}  % 缩放图表
\usepackage[table]{xcolor} % 表格颜色支持

\begin{document}
    % --- 生成封面和目录 (无需修改) ---
    \maketitle
    \makeToc
	
    % =============================================
    %                   开始正文
    % =============================================
    
    \section{Introduction to the Template}
    
    Welcome to \textbf{The University of Sydney Postgraduate Coursework Report Template}. This document is designed to provide a clean, professional, and easy-to-use LaTeX framework for your university assignments. 
    
    \subsection{How to Compile}
    This template has been optimized to run completely on \textbf{pdfLaTeX}. You no longer need to configure complex font settings or rely on XeLaTeX for basic English reports. Simply set your Overleaf compiler to pdfLaTeX, and click ``Recompile''.
    
    %-----------------------------------------------------------------------
    \section{Formatting Guide}

    \subsection{Creating Sections and Subsections}
    Organizing your report is straightforward. Use the \verb|\section{}| command for main chapters. For smaller topics, use \verb|\subsection{}| and \verb|\subsubsection{}|. The numbering is handled automatically.

    \subsubsection{A Note on Spacing}
    Do not use the \verb|\begin{spacing}{...}| environment to wrap around your section titles. The template automatically calculates the appropriate breathing room above and below your headings.

    \subsection{Citations and References}
    Academic integrity is crucial at USYD. When you need to cite a source, simply use the cite command like this: \cite{example_reference}. 
    
    Remember to update your \verb|Bibs/mybib.bib| file with your actual references. We recommend using standard English citation styles such as \verb|apalike|, \verb|harvard|, or \verb|ieeetr|.

    %-----------------------------------------------------------------------
    \section{Conclusion}
    
    You are now ready to start writing your report. Delete these instructional texts, fill in your title page information at the top of the \verb|main.tex| file, and begin typing your brilliant ideas. Good luck with your assignment!
	
    \clearpage
    % =============================================
    %                   参考文献
    % =============================================
    \bibliographystyle{apalike} % 参考文献排版风格 (可根据要求更改)
    \bibliography{Bibs/mybib}   % 指向你的 .bib 文件
    \phantomsection
    \addcontentsline{toc}{section}{References}

\end{document}